\documentclass[11pt]{article}
\usepackage[margin=1in]{geometry}
\usepackage{amsmath,amssymb,amsthm,mathtools}
\usepackage[hidelinks]{hyperref}
\newtheorem{theorem}{Theorem}
\newtheorem{lemma}[theorem]{Lemma}
\newtheorem{corollary}[theorem]{Corollary}
\newcommand{\Q}{\mathbb Q}
\newcommand{\F}{\mathbb F}
\title{Simultaneous factorial exponents, a dense-preprint power saving, and harmonic sparsity}
\author{User-supplied strengthening checked and integrated; no priority claim}
\date{17 November 2026}
\begin{document}
\maketitle
Write $F_{n,m}(T)=\sum_{j=0}^n\binom nj(mj)!T^j$ in ordinary Gamma shape one,
and let $\mathcal D$ be the set of positive integers $n$ for which some
positive integer $m\le n^{1/16}$ gives a nonsquarefree $F_{n,m}$ over $\Q$.
Put $D(X)=|\mathcal D\cap[1,X]|$ and $\varepsilon=10^{-18}$.

\begin{theorem}\label{thm:simultaneous}
There is one common exceptional set for all the indicated exponents.
\begin{enumerate}
\item Using the published 2016 Matom\"aki--Radziwi\l\l theorem,
$D(X)\ll X/(\log X)^{1/50}$ with an absolute implied constant.
\item Importing additionally the dense Theorem~9.2(ii) of the pinned
MR~II preprint, $D(X)\ll X^{1-\varepsilon}$ with an absolute implied constant.
\end{enumerate}
The second assertion uses a preprint analytic input, distinguished below
from the independent published-input first assertion. Neither numerical
constants nor starting thresholds are extracted. All-index SFC is open.
\end{theorem}

\section{A common dense irreducible-offset set}
Work on an integer dyadic block $X<n\le2X$ and put
\[
 M=\lfloor(2X)^{1/16}\rfloor,\qquad H=\lfloor X^{1/5}\rfloor,
 \qquad Y=\sqrt{2MX},\qquad Z=\sqrt{MH}.
\]
These are deterministic functions of $X$, with no separately varying $m$.
For $q_i$ the largest prime at most $i$, select
\[
 \mathcal J=\{1\le i\le H:i>(M+1)^2,\ i-q_i<\sqrt i,\ P^+(i)>Z\}.
\]
\begin{lemma}\label{lem:common}
For every $i\in\mathcal J$ and every $1\le m\le M$, $F_{i,m}$ is irreducible
in characteristic zero. Moreover
\[
 |\mathcal J|=(c_J+o(1))H,\qquad c_J=\log(32/21).
\]
This common-set statement does not apply a fixed-$m$ irreducibility-density
theorem with a growing value of $m$.
\end{lemma}
\begin{proof}
Set $p=P^+(i)>Z$. Since $p^2>MH\ge mi$ and $p\mid i$, the
\href{QUANTITATIVE_SQUAREFREENESS_ALL_EXPONENTS.tex}{general local lemma}
at offset zero gives a single edge of denominator $p$; every rational
factor has degree divisible by $p$.
Writing $i=q_i+d$, the conditions $d<\sqrt i$ and $i>(M+1)^2$ imply
$q_i>m\max(1,d)$ simultaneously for every $m\le M$.
The $q_i$-polygon has a nonhorizontal edge of width $q_i$ and denominator
$q_i$. Newton multiplication therefore supplies an irreducible rational
factor of degree at least $q_i$. Its complementary degree is at most
$d<\sqrt i<p$ and is also divisible by $p$, so is zero.

Since $Z^2=MH\ge H$, integers at most $H$ have at most one prime factor
strictly above $Z$. Disjoint prime multiples and PNT partial summation give
\[
 \#\{i\le H:P^+(i)>Z\}
 =H\sum_{Z<p\le H}\frac1p+O(\pi(H))
 =(\log(32/21)+o(1))H,
\]
using $\log H\sim(1/5)\log X$ and $\log Z\sim(21/160)\log X$.
The small-index deletion costs $O(M^2)=o(H)$.
The credited Filaseta--Rocks equation~(3.1) gives $O(H^{2/3})=o(H)$
bad preceding-gap indices, including intervals crossing the endpoint.
This proves the density of the common set, with absolute asymptotics.
\end{proof}

For $f_X(u)=\mathbf1_{\{P^+(u)\le Y\}}$, $f_X(1)=1$, disjoint large-prime
multiples are again valid because $Y^2\ge2X$. PNT gives
\[
 \mu_X(f_X):=\frac1X\sum_{X<u\le2X}f_X(u)
 =1-\log(32/17)+o(1)=:c_f+o(1),
\]
where $\log Y\sim(17/32)\log X$.
Only eventual strict margins are needed; no quantitative PNT error is assumed.
The available margin satisfies
\[
 \Delta=c_J-c_f=\log(1024/357)-1>1/20. \tag{1}
\]
For an exact certificate, the elementary positive exponential series gives
$e<2719/1000$, and $e^{1/20}<20/19$ by a geometric-series bound. Thus
\[
 e^{21/20}<\frac{2719}{1000}\frac{20}{19}<\frac{1024}{357},
\]
proving (1) without rounded transcendental numerics.

\section{Source qualification and the precise dense input}
The power-saving input is MR~II,
\href{sources/mr-ii-2007.04290v1.pdf}{arXiv:2007.04290v1, 8 July 2020},
Theorem~9.2(ii), page~44, with its Section~9 conditions, pages~42--43.
It is listed as a preprint on the author's
\href{sources/matomaki-publications-20261117.html}{publication page checked 17 November 2026}.
No published or corrected MR~II version is asserted to have been obtained.

Matom\"aki's published 2022
\href{https://doi.org/10.1112/jlms.12592}{Almost primes in almost all very short intervals}
flags an incorrect sparse-sieve calculation in the first MR~II arXiv version.
The same warning is pinned in the
\href{sources/matomaki-2022-almost-primes-accepted-v3.pdf}{accepted author version},
Section~6 footnote before Lemma~6.1; Remark~6.2 explains the missing
coprimality and resulting off-diagonal terms in the related sieve calculation.
We exclude the entire MR~II Section~10 sparse-sieve branch, including
Lemma~10.3. The selected route is the dense Theorem~9.2(ii), whose proof
on pages~45--51 uses Proposition~8.3 and earlier inputs, followed by the
separate ordinary linear sieve, Lemma~4.3, as applied in equation~(78).
At the end of the dense proof the specified replacements for part~(ii)
use Lemmas~5.4(ii), 5.5(ii) and equation~(29), rather than sparse
nonvanishing estimates. This application/dependency audit is not a
reconstruction or validation of the entire preprint. Its dense theorem
and the analytic sieve result are credited imports.
We do not use Theorem~1.7's advertised exceptional exponent.

Take exactly
\[
 \nu_1=10^{-5},\quad\nu_2=3/20,\quad\eta=1/150,
 \quad\rho=1/10,\quad\delta=1/2000,\quad J_{\rm MR}=1.
\]
Define three intervals by
\[
 P_1=X^{10^{-8}},\quad Q_1=P_2=X^{\nu_1},\quad
 Q_2=P_3=X^{\sqrt{\nu_1\nu_2}},\quad Q_3=X^{\nu_2}.
\]
All Section~9 conditions hold eventually:
$0<\nu_1<\nu_2<1/6$, $\nu_1>(\log X)^{-1/10}$,
$0<\eta<1/6-\nu_2/3$, $P_1\ge(\log X)^{2/\eta}=(\log X)^{300}$,
and $(P_1,Q_1]\subset(3/2,H]$, $2\le H\le X^{1/2}$.
The condition on $Q_J$ has a $J=1$ alternative, and spacing conditions
(50)--(51) are empty. Also
$\rho<\rho_1=1/3-2/(3\pi)$, already following from $\pi>3$, and
$0<\delta<1/1000$.

Let $\mathcal S$ require a prime divisor in each of the three intervals.
For nonnegative multiplicative $g$, the minimizing twist in the dense
statement can be chosen to be zero: for every prime,
$1-g(p)\cos(t\log p)\ge1-g(p)$.
Thus applying Theorem~9.2(ii) separately to $f_X$ and to $1$ gives
\[
 |A_x(g\mathbf1_{\mathcal S})-\mu_X(g\mathbf1_{\mathcal S})|
 \le e_X:=\frac1{2000}+\frac{10^5}{(\log X)^{1/120}},
 \quad g=f_X,1,
\]
where $A_x(g)=H^{-1}\sum_{x<u\le x+H}g(u)$, except at most
\[
 O\left(X^{1-33/10^{10}}(\log X)^2
            +X^{1-5\cdot10^{-18}}\right) \tag{2}
\]
integer starts. There are two functions, absorbed in the absolute constant;
there is no union over $m$.
Indeed the displayed Theorem~9.2(ii) exception is
$O_{\eta,\rho}(X\delta^{-2}(\log H)^2/P_1^{1/2-\nu_2-3\eta}
+X^{1-\nu_1^3/200})$.
Direct substitution gives the two savings in (2); $\rho/12=1/120$.

\section{The ordinary sieve and the single complement loss}
The linear sieve of MR~II Lemma~4.3 and its application in (78) give,
using $D=X^{\nu_2(3-1/1000)}$, $z=X^{\nu_2}$,
\[
 \mu_X(\mathcal S^c)
 \le\frac65\left(\frac1{1000}+2\sqrt{\frac1{15000}}\right)+o(1)
 <\frac{13}{625}+o(1). \tag{3}
\]
Here the three logarithmic interval ratios are
$1/1000,\sqrt{\nu_1/\nu_2},\sqrt{\nu_1/\nu_2}$.
The source uses the ordinary linear upper sieve $F(3-1/1000)<1.19$
and absorbs the remainder with $1.195<6/5$.
The level is a fixed power below $X$, so the floor remainders are $o(X)$.
The rational check $(49/6000)^2>1/15000$ proves the last bound in (3).
This argument does not use the sparse-sieve weights of Section~10.

Nonnegativity gives a sharper transfer than the generic absolute-value
triangle inequality:
\[
 \begin{aligned}
 A_x(f_X)&\le A_x(f_X\mathbf1_{\mathcal S})+1-A_x(\mathbf1_{\mathcal S})\\
 &\le\mu_X(f_X)+\mu_X(\mathcal S^c)+2e_X\\
 &\le\mu_X(f_X)+\frac{109}{5000}+o(1).
 \end{aligned}
\]
The complement occurs once. Combining this with (1) yields
\[
 \Delta-\frac{109}{5000}>
 \frac{141}{5000}>\frac1{50}. \tag{4}
\]
Choose $x=n-H-1$ for integer $n$. Then $(x,x+H]$ contains exactly
$n-H,\ldots,n-1$; discard only $O(H)$ bottom endpoint indices so that
$x\in[X,2X]$. By pigeonhole, every remaining index has at least $H/50$
distinct offsets $i\in\mathcal J$ with $P^+(n-i)>Y$.
No independence between offsets is used.

\section{Bad incidences across all exponents}
For $i\in\mathcal J$, $m\le M$, put
$R_{i,m}=\operatorname{Res}(F_{i,m},F'_{i,m})\ne0$.
The retained exact Hadamard bound gives, with an absolute constant,
\[
 \log|R_{i,m}|\le5mi^2\log(2mi),\qquad
 \sum_{m\le M}\sum_{i\in\mathcal J}\log|R_{i,m}|
 \le5M^2H^3\log(2MH). \tag{5}
\]
For example, take logarithms of
$i^i[2^i(mi)!]^{2i-1}$ and use $(mi)!\le(mi)^{mi}$.

Discard the $O(X^{2/3})$ indices with $n-q_n\ge\sqrt n$, using the pinned
Filaseta--Rocks equation~(3.1). This is one common gap condition.
For large $X$, $n>(M+1)^2$, and $Y>MH$ since
$X/(MH^2)\to\infty$. Thus every usable $p=P^+(n-i)>Y$ satisfies
$p^2>2MX\ge mn$ and $p>mH$, simultaneously for $m\le M$.
The general local transfer gives square-divisor degree
$k<\sqrt n/2<p$ and degree-preserving horizontal reduction.

If such an index fails squarefreeness for an exponent $m=m(n)\le M$,
select one failing exponent. For every one of its at least $H/50$ usable
offsets, the transfer forces $p\mid R_{i,m(n)}$.
An incidence retains the offset label even if primes coincide.
Each triple $(i,m,p)$ affects at most $X/p+1$ indices in the block,
and each $R_{i,m}$ has at most $\log|R_{i,m}|/\log Y$ large prime divisors.
Counting incidences, rather than discarding every bad triple once, gives
\[
 \begin{aligned}
 |\mathcal E_{\rm res}|
 &\ll\frac1H\frac{X/Y+1}{\log Y}
       \sum_{m\le M}\sum_{i\in\mathcal J}\log|R_{i,m}|\\
 &\ll\sqrt X\,M^{3/2}H^2\ll X^{159/160}.
 \end{aligned} \tag{6}
\]
The factor $1/H$ is essential. All constants are absolute: $M,H$ and the
analytic parameters are fixed explicit functions of $X$, and (5) retains
the full dependence on $m$. The analytic error (2) is never multiplied by $M$.

Combining (2), (6), common gaps and endpoints gives
\[
 D(2X)-D(X)\ll X^{1-33/10^{10}}(\log X)^2
 +X^{1-5\cdot10^{-18}}+X^{159/160}+X^{2/3}+X^{1/5}.
\]
Every term is $O(X^{1-10^{-18}})$, with room to absorb the logarithm.
Dyadic geometric summation proves part~(ii) of Theorem~\ref{thm:simultaneous}.
The exponent range on the whole block covers $m\le n^{1/16}$ because
$\lfloor n^{1/16}\rfloor\le M$.

\section{Independent published-2016 route}
For part~(i), retain the common $\mathcal J$, $f_X$, and incidence argument,
but use the published 2016 MR Theorem~1 directly, with $h=H$ and $\delta=1/100$.
Its short-mean error is eventually below $1/50$; (1) then again leaves
more than $H/50$ usable offsets.
Its exception is
\[
 O\left(X^{1-1/12500}(\log X)^{1/3}
                  +\frac{X}{(\log X)^{1/50}}\right).
\]
Together with (6), gaps and endpoints, dyadic summation yields part~(i),
with absolute constants and no MR~II or sparse-sieve input.
A rate for the convergence in the common density and long mean is
unnecessary: they only furnish eventual strict margins.
The 2016 result alone does not supply the reciprocal-sum consequence below.

\section{Harmonic sparsity and remaining boundary}
\begin{corollary}\label{cor:harmonic}
With the dense-preprint input of part~(ii),
\[
 \sum_{n\in\mathcal D}\frac1n<\infty.
\]
For each fixed $m\ge1$, the set of nonsquarefree $F_{n,m}$ also has a
convergent reciprocal sum.
\end{corollary}
\begin{proof}
Partial summation reduces the tail to an integral bounded by a constant
multiple of $\int_N^\infty t^{-1-\varepsilon}\,dt$, which converges.
For fixed $m$, every $n\ge m^{16}$ in its exceptional set belongs to
$\mathcal D$; the remaining initial indices are finite.
\end{proof}
The same common exceptional set controls adjacent nonzero-moment windows
for every nonzero complex $A+BY^m$ and all $m\le n^{1/16}$, by
$F_{n-1,m}=F_{n,m}-(T/n)F'_{n,m}$.
A convergent reciprocal sum does not imply finitely many exceptions or
all-index SFC. No irreducibility-density-one, arbitrary Gamma-shape,
full-preprint validation, numerical implied-constant or priority claim is made.

\href{simultaneous_squarefree.py}{Exact certificates} retain rational
parameters, error exponents and pointwise common-offset criteria.
\href{verify_simultaneous_squarefree.py}{Finite controls and source-route audit}
check these and source pins, not a numerical frequency claim.
Independent replay, formal verification and external review remain unclaimed.
\end{document}
